% !TeX root = ../MQF_Manuale_Master.tex
%2multibyte Version: 5.50.0.2960 CodePage: 65001
% MQF_Manuale_Master_SW55_Memoir.tex
% Master del manuale MQF costruito sullo shell SW5.5 "Memoir Class for Configurable Typesetting".
% Versione modulare: i capitoli scritti sono inclusi con \input.%
% POSIZIONE DEI FILE
% Questo master deve stare nella cartella 01_Manuale.
% I capitoli devono stare nella sottocartella 01_Manuale/Capitoli.
% Non mettere il master dentro Capitoli.
% I capitoli futuri restano come traccia, ma non generano pagine nel documento.
% Evita pagine bianche dovute all'apertura dei capitoli su pagina dispari.
% ------------------------------------------------------------
% Nomi italiani
% ------------------------------------------------------------
% ------------------------------------------------------------
% Macro matematiche ufficiali del progetto MQF
% Coerenti con MQF_Notazione.tex
% ------------------------------------------------------------
% ------------------------------------------------------------
% Ambienti di tipo teorema.
% I nomi interni restano compatibili con lo shell SW.
% Le etichette stampate sono in italiano.
% ------------------------------------------------------------

%TCIDATA{OutputFilter=latex2.dll}
%TCIDATA{Version=5.50.0.2960}
%TCIDATA{LaTeXparent=1,1,MQF_Manuale_Master.tex}


% Capitoli/MQF_Cap_04_Python_Probabilita_Condizionamento.tex


\chapter[Applicazioni con IA e Python]{Scenario analysis e informazione condizionata con IA e Python}

Questo capitolo mostra come gli strumenti teorici introdotti nei Capitoli
1--3 possano essere utilizzati per affrontare casi quantitativi assegnati,
trasformando definizioni e propriet\`{a} matematiche in procedure
computazionali verificabili.

Il punto di partenza \`{e} costituito dai concetti di probabilit\`{a},
variabile casuale, distribuzione, probabilit\`{a} condizionata e valore atteso
condizionato. In un problema applicativo, tuttavia, non \`{e} sufficiente
riconoscere tali oggetti in forma astratta. Occorre stabilire quali grandezze
del caso corrispondano agli oggetti teorici, quali relazioni matematiche siano
necessarie, quali quantit\`{a} debbano essere stimate e in quale ordine debbano
essere svolti i diversi passaggi.

Python viene utilizzato come strumento per rendere operative queste relazioni.
Attraverso la simulazione Monte Carlo \`{e} possibile generare scenari,
costruire distribuzioni empiriche, stimare probabilit\`{a}, momenti e quantili,
calcolare quantit\`{a} condizionate e verificare numericamente propriet\`{a}
teoriche. Il codice non costituisce quindi un fine autonomo, ma la traduzione
computazionale di una struttura teorica gi\`{a} definita.

Anche l'intelligenza artificiale viene utilizzata come strumento di supporto
alla soluzione del caso. Il suo impiego \`{e} utile soprattutto quando occorre
passare dalla formulazione del problema alla sua organizzazione operativa:
riconoscere gli oggetti teorici rilevanti, ordinare i passaggi risolutivi,
tradurre una specifica matematica in codice, individuare controlli adeguati e
verificare la coerenza degli output. L'IA non sostituisce per\`{o} il controllo
teorico e interpretativo: ogni risposta ottenuta deve essere confrontata con la
struttura del problema, con le ipotesi assegnate e con i risultati effettivamente
prodotti dal notebook.

Un passaggio particolarmente delicato consiste nel ricavare dalla traccia un
\emph{flusso logico-teorico risolutivo}. Prima di procedere al codice occorre
individuare le grandezze assegnate, la variabile o le variabili di interesse,
gli eventi rilevanti, l'informazione disponibile, le relazioni funzionali, le
quantit\`{a} teoriche da calcolare e le propriet\`{a} che consentono di
collegarle. Il flusso deve inoltre rispettare un ordine logico: ogni passaggio
deve utilizzare oggetti gi\`{a} definiti e preparare il passaggio successivo.

Una volta definito il flusso teorico, il problema deve essere tradotto in una
\emph{scomposizione operativa del notebook in tappe input--output}. Ogni tappa
deve specificare quali input vengono utilizzati, quale operazione viene
eseguita, quale output deve essere prodotto e quale controllo deve consentire
di valutarne la coerenza. In questo modo il notebook non diventa una semplice
sequenza di istruzioni Python, ma una rappresentazione osservabile del
ragionamento quantitativo.

Il caso sviluppato in questo capitolo riguarda la perdita di un portafoglio
obbligazionario a un periodo in presenza di differenti regimi informativi del
mercato dei tassi. La struttura del problema consente di utilizzare
congiuntamente una partizione dello spazio degli eventi, distribuzioni
condizionate, una variabile casuale di perdita, probabilit\`{a} di superamento
di soglia, valori attesi condizionati e formula del valore atteso totale. La
simulazione consente di osservare come queste quantit\`{a} cambino al variare
dell'informazione disponibile e di confrontare sistematicamente formulazione
teorica e stima empirica.

Accanto al caso sviluppato nel capitolo viene considerato un secondo caso, da
affrontare in autonomia. Il contesto economico cambia, ma resta invariato il
metodo di lavoro: riconoscimento degli oggetti teorici, costruzione del flusso
risolutivo, scomposizione del notebook, formulazione dei prompt, implementazione
in Python, verifica degli output e interpretazione finale.

L'obiettivo del capitolo \`{e} quindi duplice. Da un lato, mostrare come i
concetti probabilistici e condizionati dei capitoli precedenti vengano
utilizzati nella soluzione di un caso quantitativo. Dall'altro, sviluppare un
metodo di lavoro riutilizzabile, nel quale teoria finanziaria, Python e
intelligenza artificiale contribuiscono a funzioni diverse ma coordinate della
stessa procedura risolutiva.

\section{Obiettivi del capitolo e raccordo con i Capitoli 1--3}

I primi tre capitoli hanno costruito progressivamente il linguaggio necessario
per rappresentare l'incertezza finanziaria. Il Capitolo 1 ha introdotto eventi,
probabilit\`{a}, partizioni e condizionamento; il Capitolo 2 ha trasformato gli
esiti aleatori in grandezze numeriche mediante le variabili casuali e le loro
distribuzioni; il Capitolo 3 ha mostrato come la valutazione di tali grandezze
cambi quando si dispone di informazione parziale.

In questo capitolo questi concetti vengono utilizzati congiuntamente per
risolvere casi quantitativi mediante simulazione. L'obiettivo non \`{e}
introdurre nuovi risultati probabilistici, ma imparare a riconoscere quali
strumenti teorici siano necessari in un problema assegnato e come possano
essere tradotti in una procedura computazionale controllabile.

Il passaggio fondamentale pu\`{o} essere rappresentato come segue:

\[
\begin{array}{c}
\text{problema finanziario}
\\[4pt]
\downarrow
\\[4pt]
\text{oggetti teorici rilevanti}
\\[4pt]
\downarrow
\\[4pt]
\text{quantit\`{a} da calcolare o stimare}
\\[4pt]
\downarrow
\\[4pt]
\text{procedura computazionale}
\\[4pt]
\downarrow
\\[4pt]
\text{output e controlli}.
\end{array}
\]

Python interviene nella traduzione computazionale degli oggetti gi\`{a}
identificati sul piano teorico. L'intelligenza artificiale pu\`{o} assistere
questo passaggio, ma il suo utilizzo rimane subordinato alla corretta
specificazione del problema e alla verifica dei risultati prodotti.

\subsection{Dagli eventi agli scenari informativi}

Nel Capitolo 1 una partizione
\[
\{A_{1},\ldots,A_{m}\}
\]
dello spazio campionario \`{e} stata introdotta come una famiglia di eventi
disgiunti la cui unione coincide con $\Omega$. Nei casi applicativi, i blocchi
della partizione possono rappresentare differenti condizioni di mercato o
stati informativi.

Deve quindi valere
\[
\Omega=A_{1}\cup\cdots\cup A_{m},
\qquad
A_{i}\cap A_{j}=\emptyset
\quad\text{per }i\neq j,
\]
e
\[
\sum_{g=1}^{m}\Prob(A_{g})=1.
\]

In una simulazione, la probabilit\`{a} teorica $\Prob(A_g)$ trova una
controparte empirica nella frequenza con cui il regime $A_g$ compare nel
campione. Se vengono generate $M$ simulazioni e $M_g$ appartengono ad $A_g$,
si ottiene
\[
\widehat{\Prob}(A_g)
=
\frac{M_g}{M}.
\]

Il confronto tra probabilit\`{a} teoriche e frequenze empiriche costituisce un
primo controllo della simulazione.

\subsection{Dalle variabili casuali alle distribuzioni empiriche}

Il Capitolo 2 ha introdotto una variabile casuale come funzione
\[
X:\Omega\longrightarrow\mathbb{R}.
\]

Nei casi quantitativi, $X$ pu\`{o} rappresentare un rendimento, una variazione
di tasso, una perdita, un costo o un payoff. La distribuzione teorica di $X$
viene resa osservabile mediante un campione Monte Carlo
\[
X^{(1)},X^{(2)},\ldots,X^{(M)}.
\]

Da tale campione possono essere ricavate stime delle quantit\`{a} teoriche
studiate nei capitoli precedenti. Ad esempio,
\[
\widehat{\E}[X]
=
\frac{1}{M}\sum_{k=1}^{M}X^{(k)}
\]
\`{e} la media campionaria, mentre
\[
\widehat{\Prob}(X>x)
=
\frac{1}{M}
\sum_{k=1}^{M}
\1_{\{X^{(k)}>x\}}
\]
stima una probabilit\`{a} di superamento di soglia.

La distribuzione empirica permette inoltre di calcolare deviazione standard,
quantili e altre statistiche descrittive. La simulazione non modifica quindi
l'oggetto teorico: ne costruisce una rappresentazione empirica sulla quale
possono essere effettuati calcoli e controlli.

\subsection{Dal condizionamento teorico alle stime per gruppi informativi}

Il Capitolo 3 ha mostrato che
\[
\E[X\mid A_g]
\]
rappresenta la valutazione media di $X$ quando \`{e} noto che l'evento $A_g$
si \`{e} verificato.

Nel campione simulato la stessa idea viene tradotta selezionando le osservazioni
appartenenti al regime considerato. Se
\[
I_g=\{k:\omega_k\in A_g\},
\qquad
M_g=\# I_g,
\]
la corrispondente stima \`{e}
\[
\widehat{\E}[X\mid A_g]
=
\frac{1}{M_g}
\sum_{k\in I_g}X^{(k)}.
\]

Se
\[
\G=\sigma(A_1,\ldots,A_m),
\]
il valore atteso condizionato rispetto all'informazione generata dalla
partizione assume la forma
\[
\E[X\mid\G]
=
\sum_{g=1}^{m}
\E[X\mid A_g]\1_{A_g}.
\]

La formula del valore atteso totale fornisce inoltre un controllo centrale:
\[
\E[X]
=
\sum_{g=1}^{m}
\Prob(A_g)\E[X\mid A_g].
\]

Nel campione simulato la corrispondente ricomposizione \`{e}
\[
\widehat{\E}[X]
=
\sum_{g=1}^{m}
\widehat{\Prob}(A_g)
\widehat{\E}[X\mid A_g].
\]

Una propriet\`{a} teorica pu\`{o} quindi svolgere una doppia funzione:
giustificare un calcolo e fornire un controllo della sua implementazione.

\subsection{Obiettivi operativi}

Al termine del capitolo dovr\`{a} risultare possibile:

\begin{enumerate}
\item riconoscere in un caso assegnato eventi, variabili casuali, stati
informativi e quantit\`{a} condizionate introdotti nei Capitoli 1--3;

\item distinguere una quantit\`{a} teorica dalla corrispondente stima empirica
ottenuta mediante simulazione;

\item individuare le relazioni matematiche che collegano gli input del problema
alla variabile finale di interesse;

\item organizzare tali relazioni in un ordine logico prima della scrittura del
codice;

\item trasformare il percorso teorico in una procedura Python con output e
controlli espliciti;

\item utilizzare l'IA come supporto alla ricognizione, alla traduzione
computazionale e alla verifica;

\item valutare la coerenza dei risultati numerici rispetto alle ipotesi e alle
propriet\`{a} teoriche utilizzate.
\end{enumerate}

Il raccordo tra teoria e applicazione non consiste quindi nel sostituire le
formule con il codice. Consiste nel riconoscere la stessa struttura
quantitativa prima come modello probabilistico e successivamente come procedura
computazionale verificabile.

\section{Il metodo di lavoro nei capitoli applicativi}

Nei capitoli applicativi il problema quantitativo non viene affrontato
scrivendo immediatamente il codice. La procedura parte dalla comprensione della
traccia, passa attraverso la ricostruzione del modello teorico e conduce
successivamente alla costruzione del notebook Jupyter.

L'obiettivo \`{e} mantenere riconoscibile la sequenza

\[
\begin{array}{c}
\text{problema finanziario}
\\[4pt]
\downarrow
\\[4pt]
\text{modello teorico}
\\[4pt]
\downarrow
\\[4pt]
\text{procedura computazionale}
\\[4pt]
\downarrow
\\[4pt]
\text{output}
\\[4pt]
\downarrow
\\[4pt]
\text{controllo e interpretazione}.
\end{array}
\]

L'intelligenza artificiale pu\`{o} assistere questo percorso, ma le diverse
funzioni svolte dallo strumento devono rimanere riconoscibili e verificabili.

\subsection{La Scheda Caso come specifica del problema}

Ogni applicazione prende avvio da una Scheda Caso, che definisce il perimetro
del lavoro. In particolare, stabilisce:

\begin{itemize}
\item il contesto e la domanda quantitativa;

\item le grandezze, le variabili e le ipotesi del modello;

\item i parametri assegnati;

\item le quantit\`{a} da calcolare o stimare;

\item gli output e i controlli richiesti;

\item i limiti entro i quali interpretare i risultati.
\end{itemize}

La Scheda Caso non contiene la soluzione. Stabilisce invece gli elementi che
non devono essere modificati durante lo svolgimento.

Anche nell'interazione con l'IA questo vincolo rimane essenziale. Una risposta
che modifica autonomamente una variabile, una distribuzione, una soglia,
un'ipotesi o il significato finanziario del problema devia dalla specifica
assegnata.

I diversi strumenti svolgono quindi funzioni distinte:

\begin{center}
\begin{tabular}{p{4.0cm}p{7.4cm}}
Scheda Caso & definisce che cosa deve essere risolto; \\
Interazione con l'IA & assiste l'organizzazione e la traduzione operativa; \\
Notebook Jupyter & documenta la procedura implementata e i risultati; \\
Controlli & verificano la coerenza tra specifica, codice e output.
\end{tabular}
\end{center}

\subsection{Caso sviluppato e caso autonomo}

Il primo caso viene sviluppato nel capitolo e rende osservabile l'intera
sequenza di lavoro: ricognizione teorica, organizzazione dei passaggi,
costruzione del notebook, controllo degli output e interpretazione.

Il Caso TakeHome richiede invece il trasferimento autonomo dello stesso metodo
a un problema differente. Non consiste nel riprodurre il notebook gi\`{a}
sviluppato cambiandone parametri o denominazioni. Devono essere nuovamente
riconosciuti gli oggetti teorici rilevanti e ricostruita la procedura
risolutiva.

Cambia quindi il contenuto del problema; rimane invariato il metodo di lavoro.

\subsection{Tre funzioni dell'interazione con l'IA}

L'interazione con l'IA viene distinta in tre regimi.

\begin{center}
\small
\begin{tabular}{p{1.7cm}p{2.9cm}p{6.1cm}}
Regime & Funzione prevalente & Attivit\`{a} principale \\
\hline
A & Ricognizione teorico-modellistica &
riconoscere gli oggetti teorici, verificare relazioni e ordinare il
ragionamento; \\

B & Traduzione operativa &
trasformare una specifica teorica validata in celle Markdown, codice Python,
tabelle, grafici e output; \\

C & Verifica critica &
controllare errori, anomalie, incoerenze o difformit\`{a} rispetto alla
Scheda Caso.
\end{tabular}
\end{center}

Nel Regime A viene chiarito che cosa deve essere calcolato e perch\'{e}; nel
Regime B una specifica gi\`{a} validata viene resa computabile; nel Regime C
viene verificata una criticit\`{a} individuata nel lavoro gi\`{a} prodotto.

\subsection{La sequenza dei prompt}

L'interazione viene organizzata mediante prompt con funzioni differenti:

\[
\begin{array}{c}
\text{Prompt zero}
\\
\downarrow
\\
\text{Prompt 1}
\\
\downarrow
\\
\text{Prompt 2}
\\
\downarrow
\\
\text{Prompt 3}
\\
\downarrow
\\
\text{prompt di tappa}
\\
\downarrow
\\
\text{eventuali verifiche in Regime C}.
\end{array}
\]

Il \textbf{Prompt zero} definisce il contesto generale e le regole
dell'interazione.

Il \textbf{Prompt 1} acquisisce la Scheda Caso come specifica vincolante e
produce l'inquadramento iniziale del notebook.

Il \textbf{Prompt 2} riguarda il flusso logico-teorico risolutivo. Deve partire
da una proposta iniziale che renda visibili gli elementi teorici individuati
prima dell'interazione con l'IA.

Il \textbf{Prompt 3} trasforma il flusso validato in una scomposizione
operativa del notebook. Anche in questo caso deve esistere una proposta
iniziale di tappe.

I \textbf{prompt di tappa} intervengono soltanto dopo la definizione del
percorso teorico e della struttura operativa. Le richieste in
\textbf{Regime C} vengono invece utilizzate quando emerge una criticit\`{a}
specifica.

La sequenza impedisce che la produzione del codice preceda la comprensione e
l'organizzazione del problema.

\subsection{Il notebook come rappresentazione del ragionamento}

Il notebook non deve essere una successione di celle di codice prive di
raccordo. Ogni tappa deve permettere di riconoscere:

\begin{enumerate}
\item l'obiettivo dell'operazione;

\item il codice che la realizza;

\item l'output prodotto;

\item il controllo effettuato;

\item quando necessario, il relativo commento interpretativo.
\end{enumerate}

Il codice deve quindi rispondere a una domanda gi\`{a} formulata; tabelle e
grafici devono rappresentare quantit\`{a} precedentemente individuate; i
controlli devono verificare relazioni con un significato teorico o
computazionale preciso.

\subsection{Il tracciato dell'interazione}

Quando richiesto, l'interazione con l'IA viene documentata mediante un
tracciato separato dal notebook. La sua funzione \`{e} rendere osservabile il
modo in cui l'assistenza \`{e} stata utilizzata.

Nel tracciato devono risultare riconoscibili la richiesta formulata, il regime
utilizzato, il contributo iniziale, la risposta ricevuta, la decisione assunta
e i controlli successivamente effettuati.

Una risposta dell'IA non viene trasferita automaticamente nel notebook. Pu\`{o}
contenere errori, modificare implicitamente il modello, confondere quantit\`{a}
teoriche e stime empiriche oppure produrre codice non coerente con la
specifica.

Il criterio generale \`{e} quindi

\[
\begin{array}{c}
\text{proposta}
\\[3pt]
\downarrow
\\[3pt]
\text{risposta IA}
\\[3pt]
\downarrow
\\[3pt]
\text{verifica}
\\[3pt]
\downarrow
\\[3pt]
\text{decisione}.
\end{array}
\]

L'uso dell'IA diventa formativo quando la risposta ricevuta viene trattata
come una proposta da controllare, non come una soluzione da assumere
automaticamente.


\section{Dal Caso Aula al flusso logico-teorico risolutivo}

Il primo passaggio realmente autonomo nella soluzione di un caso consiste nel
ricostruire la struttura teorica del problema prima di procedere alla
programmazione.

La Scheda Caso fornisce variabili, parametri, ipotesi, output richiesti e
controlli. Non fornisce per\`{o} l'ordine logico con cui tali elementi devono
essere utilizzati. Questo ordine deve essere ricostruito.

Nel caso sviluppato in questo capitolo si considera un portafoglio
obbligazionario a tasso fisso esposto a uno shock del rendimento di mercato.
L'informazione disponibile distingue tre regimi:
\[
A_{1},\qquad A_{2},\qquad A_{3},
\]
che formano una partizione dello spazio degli eventi. Condizionatamente al
regime realizzato, lo shock di rendimento $\Delta y$ segue una distribuzione
normale con parametri specifici:
\[
\Delta y\mid A_{g}
\sim
\mathcal{N}(\mu_{g},\sigma_{g}^{2}).
\]

La variazione del valore del portafoglio viene approssimata mediante la
duration modificata. Indicando con $V_{0}$ il valore iniziale del portafoglio
e con $D$ la duration modificata, la perdita viene definita come
\[
L=D V_{0}\Delta y.
\]

La quantit\`{a} $L$ costituisce la variabile centrale del caso. Gli output
richiesti riguardano la sua distribuzione empirica, la probabilit\`{a} di
superamento di una soglia, le medie condizionate rispetto ai regimi
informativi e la ricomposizione della media complessiva.

Queste informazioni sono sufficienti per definire il problema, ma non ancora
per stabilire come risolverlo. Occorre trasformarle in una sequenza teorica
ordinata.

\subsection{La funzione del Prompt 2}

Il Prompt 2 ha lo scopo di costruire il \emph{flusso logico-teorico
risolutivo}. La sua parte variabile deve essere preparata prima dell'interazione
con l'IA.

Non \`{e} quindi sufficiente formulare una richiesta del tipo:

\begin{quote}
\small
\noindent{\ttfamily\raggedright
Costruisci il flusso logico-teorico risolutivo del caso.\par}
\end{quote}

Una richiesta di questo tipo trasferirebbe all'IA il compito di individuare
gli elementi teorici rilevanti e di scegliere la struttura della soluzione.

La parte variabile del Prompt 2 deve invece contenere una prima ricognizione
proposta autonomamente. Tale ricognizione non deve essere necessariamente
completa o perfettamente ordinata. Deve per\`{o} mostrare che sono stati
riconosciuti gli oggetti fondamentali del problema.

L'IA interviene successivamente per verificare, completare e ordinare la
proposta.

\subsection{Come leggere una Scheda Caso prima di formulare il Prompt 2}

Una lettura efficace della Scheda Caso pu\`{o} essere organizzata attraverso
alcune domande successive. Le domande non costituiscono una formula da
applicare meccanicamente, ma una guida per separare gli elementi che svolgono
funzioni differenti nella soluzione.

\paragraph{1. Qual \`{e} la variabile finale di interesse?}

Occorre innanzitutto individuare la grandezza sulla quale si concentrano gli
output e l'interpretazione.

Nel caso considerato la variabile centrale \`{e}
\[
L=D V_{0}\Delta y.
\]

Questa identificazione orienta l'intera ricognizione. Le altre grandezze
devono essere considerate in funzione del ruolo che svolgono nella
determinazione o nell'analisi di $L$.

\paragraph{2. Da quali grandezze dipende la variabile finale?}

Una volta identificata la variabile centrale, occorre risalire alle grandezze
che la generano.

Nel caso in esame:
\[
\Delta y
\longrightarrow
L.
\]

La relazione
\[
L=D V_{0}\Delta y
\]
mostra che $V_{0}$ e $D$ sono parametri assegnati, mentre $\Delta y$ \`{e} la
variabile casuale che genera l'incertezza della perdita.

Questo passaggio permette di distinguere tra:

\begin{itemize}
\item parametri deterministici;
\item variabili casuali;
\item variabili ottenute mediante trasformazione di altre variabili.
\end{itemize}

\paragraph{3. Da dove deriva l'incertezza della variabile di input?}

La distribuzione di $\Delta y$ non \`{e} unica. Dipende dal regime
informativo realizzato.

Occorre quindi riconoscere la struttura
\[
A_{g}
\longrightarrow
\Delta y\mid A_{g}
\longrightarrow
L.
\]

Gli eventi $A_{1},A_{2},A_{3}$ non rappresentano semplicemente etichette
descrittive. Determinano quale distribuzione condizionata viene utilizzata
nella generazione dello shock.

Il richiamo teorico rilevante appartiene quindi al Capitolo 1, per la
partizione degli eventi, e al Capitolo 3, per il ruolo dell'informazione
condizionante.

\paragraph{4. Quali quantit\`{a} devono essere calcolate sulla variabile finale?}

Dalla Scheda Caso devono essere estratte le domande quantitative.

Nel caso considerato esse comprendono, tra le altre:
\[
\E[L],
\qquad
\Prob(L>\ell),
\qquad
\E[L\mid A_{g}],
\qquad
\E[L\mid\G].
\]

Questa fase \`{e} importante perch\'{e} impedisce di confondere la variabile
centrale con le quantit\`{a} che devono essere calcolate su di essa.

$L$ \`{e} una variabile casuale; $\E[L]$ \`{e} una sua caratteristica media;
$\Prob(L>\ell)$ \`{e} una probabilit\`{a} di coda;
$\E[L\mid A_{g}]$ \`{e} una media condizionata.

\paragraph{5. Quali propriet\`{a} teoriche collegano le quantit\`{a} richieste?}

Non tutti gli output devono essere trattati come calcoli indipendenti.

La partizione informativa
\[
\G=\sigma(A_{1},A_{2},A_{3})
\]
consente di scrivere
\[
\E[L\mid\G]
=
\sum_{g=1}^{3}
\E[L\mid A_{g}]\1_{A_{g}}.
\]

La formula del valore atteso totale permette inoltre di ricomporre
\[
\E[L]
=
\sum_{g=1}^{3}
\Prob(A_{g})\E[L\mid A_{g}].
\]

Questa relazione non produce soltanto una quantit\`{a} aggiuntiva. Fornisce
anche un controllo teorico sulla coerenza della soluzione.

\paragraph{6. Quali risultati teorici possono diventare controlli?}

Durante la ricognizione conviene distinguere tra formule utilizzate per
ottenere un risultato e formule utilizzate per verificarlo.

Ad esempio:
\[
\sum_{g=1}^{3}\Prob(A_{g})=1
\]
\`{e} una propriet\`{a} della partizione probabilistica;

\[
L>\ell
\quad\Longleftrightarrow\quad
\Delta y>\frac{\ell}{D V_{0}}
\]
fornisce un'equivalenza tra due rappresentazioni dello stesso evento;

\[
\E[L]
=
\E\left[\E[L\mid\G]\right]
\]
fornisce una condizione di ricomposizione.

Tutte queste relazioni possono diventare controlli del successivo
procedimento computazionale.

\subsection{Costruire una catena di dipendenze}

Le risposte alle domande precedenti devono essere trasformate in una sequenza.

Un modo utile per farlo consiste nel chiedersi, per ogni elemento individuato:

\begin{quote}
\small
\noindent{\ttfamily\raggedright
Questo oggetto deve essere conosciuto prima o dopo gli altri?
Da quale informazione dipende?
Quale quantit\`{a} permette di costruire?
\par}
\end{quote}

Nel caso considerato emerge una catena essenziale:
\[
\{A_{1},A_{2},A_{3}\}
\longrightarrow
\Delta y\mid A_{g}
\longrightarrow
\Delta y
\longrightarrow
L
\longrightarrow
\left\{
\begin{array}{c}
\Prob(L>\ell),\\
\E[L],\\
\E[L\mid A_{g}]
\end{array}
\right\}
\longrightarrow
\E[L\mid\G]
\longrightarrow
\E[\E[L\mid\G]].
\]

Questa rappresentazione non costituisce ancora il flusso definitivo, ma
consente di controllare che non siano stati invertiti rapporti di dipendenza.

Non sarebbe, ad esempio, possibile analizzare la distribuzione della perdita
prima di avere definito come viene generato $\Delta y$, oppure costruire
$\E[L\mid\G]$ prima di avere determinato i valori condizionati
$\E[L\mid A_{g}]$.

\subsection{Teoria e controparte Monte Carlo}

Un'ulteriore distinzione deve essere effettuata prima della formulazione del
Prompt 2: quella tra quantit\`{a} teorica e quantit\`{a} empirica.

La quantit\`{a}
\[
\E[L\mid A_{g}]
\]
\`{e} un oggetto teorico.

Nel campione Monte Carlo, la sua controparte \`{e}
\[
\widehat{\E}[L\mid A_{g}]
=
\frac{1}{M_{g}}
\sum_{k\in I_{g}}L^{(k)}.
\]

Analogamente,
\[
\Prob(L>\ell)
\]
ha come controparte empirica
\[
\widehat{\Prob}(L>\ell)
=
\frac{1}{M}
\sum_{k=1}^{M}
\1_{\{L^{(k)}>\ell\}}.
\]

La distinzione \`{e} essenziale. Il flusso logico-teorico deve identificare
l'oggetto matematico che interessa; il notebook dovr\`{a} successivamente
costruirne la stima empirica.

Confondere i due livelli pu\`{o} produrre formulazioni improprie. Ad esempio,
una definizione analitica mediante integrale pu\`{o} essere corretta sul piano
teorico ma poco informativa quando l'obiettivo della tappa \`{e} specificare
come una quantit\`{a} verr\`{a} stimata nel campione simulato.

\subsection{Preparare la parte variabile del Prompt 2}

Dopo la ricognizione, la parte variabile del Prompt 2 pu\`{o} essere formulata
come una proposta sintetica.

Non \`{e} necessario anticipare una soluzione completa. Una proposta iniziale
pu\`{o} assumere, ad esempio, la forma:

\begin{quote}
\small
\noindent{\ttfamily\raggedright
Gli elementi che sembrano necessari per costruire il flusso risolutivo sono:
la partizione dei tre regimi informativi; la distribuzione dello shock
condizionata al regime; la relazione tra shock e perdita; il superamento della
soglia di perdita; il calcolo delle medie condizionate nei diversi regimi; la
costruzione del valore atteso condizionato rispetto all'informazione generata
dalla partizione; la verifica finale mediante il valore atteso totale.

Verifica se questi elementi sono sufficienti, se l'ordine \`{e} corretto e se
manca qualche passaggio teorico necessario. Non modificare la Scheda Caso e
non produrre ancora codice Python.
\par}
\end{quote}

La qualit\`{a} della proposta non dipende dal fatto che tutti gli elementi
siano gi\`{a} collocati nella posizione definitiva. La funzione
dell'interazione in questa fase \`{e} proprio verificare e ordinare la
ricognizione.

Deve invece essere riconoscibile il contributo autonomo iniziale: le
principali componenti teoriche devono essere state individuate prima della
risposta dell'IA.

\subsection{Dalla proposta iniziale al flusso validato}

La risposta dell'IA deve essere valutata sulla base di alcuni criteri.

Occorre verificare che:

\begin{enumerate}
\item tutti i passaggi derivino dalla Scheda Caso;

\item non vengano introdotte variabili o ipotesi non previste;

\item l'ordine rispetti le dipendenze logiche tra gli oggetti;

\item siano distinte quantit\`{a} teoriche e stime empiriche;

\item ogni formula abbia una funzione effettiva nella soluzione;

\item i controlli proposti siano collegati a propriet\`{a} verificabili;

\item non venga anticipato codice Python.
\end{enumerate}

Nel caso considerato, il flusso validato pu\`{o} essere sintetizzato come segue.

\begin{center}
\small
\begin{tabular}{p{0.8cm}p{2.7cm}p{3.6cm}p{3.5cm}p{2.4cm}}
\textbf{Passo} &
\textbf{Finalit\`{a}} &
\textbf{Oggetto teorico} &
\textbf{Applicazione nel caso} &
\textbf{Output o controllo}
\\ \hline

1 &
Rappresentare l'informazione &
Partizione
$\{A_{1},A_{2},A_{3}\}$ &
Identificazione dei tre regimi del mercato dei tassi &
Somma delle probabilit\`{a} pari a uno
\\

2 &
Specificare l'incertezza condizionata &
$\Delta y\mid A_{g}\sim
\mathcal{N}(\mu_{g},\sigma_{g}^{2})$ &
Distribuzione dello shock diversa nei tre regimi &
Confronto tra parametri teorici e campionari
\\

3 &
Generare lo shock non condizionato &
Combinazione delle distribuzioni condizionate secondo
$\Prob(A_{g})$ &
Estrazione del regime e successiva estrazione di $\Delta y$ &
Frequenze empiriche dei regimi
\\

4 &
Costruire la perdita &
$L=D V_{0}\Delta y$ &
Trasformazione dello shock in perdita monetaria &
Coerenza della trasformazione e dei segni
\\

5 &
Analizzare il rischio di soglia &
$\Prob(L>\ell)$ e
$L>\ell\Longleftrightarrow
\Delta y>\ell/(D V_{0})$ &
Probabilit\`{a} di perdita superiore alla soglia assegnata &
Equivalenza dei due eventi
\\

6 &
Valutare la perdita nei diversi regimi &
$\E[L\mid A_{g}]$ &
Media della perdita condizionata al regime osservato &
Confronto tra i tre valori condizionati
\\

7 &
Rappresentare la valutazione coerente con l'informazione &
$\E[L\mid\G]$ &
Valore costante su ciascun blocco della partizione &
Tre valori associati ai tre regimi
\\

8 &
Ricomporre la valutazione complessiva &
$\E[L]=\E[\E[L\mid\G]]$ &
Confronto tra media globale e media delle medie condizionate &
Verifica del valore atteso totale
\end{tabular}
\end{center}

La tabella non deve essere interpretata come uno schema da memorizzare. Essa
\`{e} il risultato della ricognizione svolta sullo specifico problema.

In un caso differente potranno cambiare variabili, formule, propriet\`{a} e
numero dei passaggi. Rimane invece invariato il metodo utilizzato per
ricavarli:

\[
\text{variabile finale}
\longrightarrow
\text{dipendenze}
\longrightarrow
\text{informazione}
\longrightarrow
\text{quantit\`{a} richieste}
\longrightarrow
\text{propriet\`{a} teoriche}
\longrightarrow
\text{controlli}
\longrightarrow
\text{ordine logico}.
\]

Il Prompt 2 diventa cos\`{\i} il punto di arrivo di una ricognizione teorica,
non il suo punto di partenza.


\section{Dal flusso teorico alle tappe del notebook}

Il flusso logico-teorico risolutivo stabilisce quali oggetti devono essere
utilizzati e in quale ordine logico. Non stabilisce ancora come il lavoro debba
essere organizzato nel notebook Jupyter.

Il passaggio successivo consiste quindi nel trasformare il flusso teorico in
una successione di tappe operative. Ogni tappa deve costituire un'unit\`{a} di
lavoro sufficientemente autonoma, ma deve anche collegarsi alle tappe
precedenti e successive.

La struttura generale \`{e}

\[
\text{flusso logico-teorico}
\longrightarrow
\text{tappe operative}
\longrightarrow
\text{notebook Jupyter}.
\]

Il Prompt 3 viene utilizzato per verificare questa trasformazione. Anche in
questo caso, tuttavia, la parte variabile del prompt deve contenere una prima
proposta elaborata prima dell'interazione con l'IA.

\subsection{La funzione del Prompt 3}

Il Prompt 3 non deve chiedere all'IA di costruire liberamente il notebook.

Una richiesta del tipo

\begin{quote}
\small
\noindent{\ttfamily\raggedright
Dividi il problema in tappe e costruisci il notebook.
\par}
\end{quote}

sarebbe troppo generica. L'IA dovrebbe decidere autonomamente quali operazioni
raggruppare, quali output produrre, quali controlli inserire e quale ordine
computazionale seguire.

La funzione del Prompt 3 \`{e} invece pi\`{u} delimitata: partire dal flusso
teorico gi\`{a} validato e verificare una proposta iniziale di scomposizione
operativa.

La parte variabile deve quindi rendere visibile una prima risposta alla
domanda:

\begin{quote}
\small
\noindent{\ttfamily\raggedright
Come pu\`{o} essere trasformato il flusso teorico in una successione di
blocchi computazionali collegati tra loro?
\par}
\end{quote}

\subsection{Dai passi teorici alle operazioni}

Non esiste necessariamente una corrispondenza uno-a-uno tra passi teorici e
tappe del notebook.

Alcuni passaggi teorici possono essere realizzati nella stessa tappa
computazionale. Al contrario, un singolo oggetto teorico pu\`{o} richiedere
pi\`{u} operazioni quando la sua costruzione comporta simulazione, produzione
di output e controllo.

Per decidere come raggruppare i passaggi conviene esaminare ciascun elemento
del flusso attraverso quattro domande:

\begin{enumerate}
\item quali informazioni devono essere gi\`{a} disponibili?

\item quale operazione deve essere eseguita?

\item quale risultato osservabile deve essere prodotto?

\item quale controllo deve essere effettuato prima di procedere?
\end{enumerate}

Queste quattro domande definiscono la struttura elementare di una tappa:

\[
\boxed{
\text{input}
\longrightarrow
\text{operazione}
\longrightarrow
\text{output}
\longrightarrow
\text{controllo}
}
\]

L'output prodotto pu\`{o} inoltre diventare un input della tappa successiva.
Per questa ragione la scomposizione deve essere considerata come una catena di
dipendenze, non come un semplice elenco di attivit\`{a}.

\subsection{Individuare gli input}

Gli input di una tappa comprendono tutto ci\`{o} che deve essere disponibile
prima che l'operazione possa essere eseguita.

Possono essere:

\begin{itemize}
\item parametri assegnati dalla Scheda Caso;

\item distribuzioni teoriche gi\`{a} specificate;

\item variabili simulate in una tappa precedente;

\item tabelle o vettori prodotti precedentemente;

\item soglie o altri valori necessari al calcolo.
\end{itemize}

Nel caso sviluppato, ad esempio, la simulazione dello shock di rendimento non
pu\`{o} essere effettuata senza avere prima definito:

\[
\Prob(A_{g}),\qquad
\mu_{g},\qquad
\sigma_{g},
\]
per ciascun regime.

Analogamente, la perdita
\[
L=D V_{0}\Delta y
\]
non pu\`{o} essere calcolata prima di disporre dello shock simulato
$\Delta y$ e dei parametri $D$ e $V_{0}$.

L'identificazione degli input impedisce quindi di collocare una tappa in una
posizione incompatibile con le sue dipendenze.

\subsection{Definire l'operazione}

L'operazione descrive che cosa deve essere fatto utilizzando gli input
disponibili.

La descrizione deve rimanere indipendente dalla sintassi Python. In questa fase
non interessa ancora stabilire quale funzione o quale libreria verr\`{a}
utilizzata. Interessa stabilire il significato quantitativo dell'operazione.

Sono esempi di operazioni:

\begin{itemize}
\item simulare il regime informativo secondo le probabilit\`{a} assegnate;

\item simulare lo shock condizionatamente al regime;

\item trasformare lo shock nella perdita monetaria;

\item calcolare statistiche descrittive della perdita;

\item identificare le osservazioni che superano una soglia;

\item raggruppare le perdite per regime informativo;

\item calcolare medie e probabilit\`{a} condizionate;

\item ricomporre la media complessiva mediante le frequenze dei regimi.
\end{itemize}

Una buona descrizione dell'operazione risponde quindi alla domanda
\emph{che cosa deve essere calcolato?}, non ancora alla domanda
\emph{con quale istruzione Python deve essere calcolato?}.

\subsection{Definire l'output}

Ogni tappa deve produrre almeno un risultato osservabile.

Un output pu\`{o} assumere forme differenti:

\begin{itemize}
\item una variabile simulata;

\item un vettore di osservazioni;

\item una tabella;

\item una statistica;

\item una probabilit\`{a} stimata;

\item un grafico;

\item una quantit\`{a} che verr\`{a} utilizzata nella tappa successiva.
\end{itemize}

La presenza di un output esplicito svolge una funzione importante. Permette di
stabilire se la tappa ha realmente prodotto ci\`{o} che il flusso teorico
richiedeva.

Ad esempio, la tappa dedicata alla simulazione dei regimi non dovrebbe
limitarsi a creare un vettore di etichette. Deve rendere osservabile almeno la
distribuzione delle frequenze simulate, perch\'{e} queste ultime possono essere
confrontate con le probabilit\`{a} teoriche assegnate.

\subsection{Associare un controllo a ogni tappa}

Il controllo non deve essere aggiunto soltanto alla fine del notebook.

Ogni tappa dovrebbe concludersi, quando possibile, con una verifica coerente
con la natura dell'operazione svolta.

Nel caso considerato possono essere utilizzati, ad esempio:

\[
\sum_{g=1}^{3}\widehat{\Prob}(A_{g})=1,
\]

il confronto tra
\[
\widehat{\Prob}(A_{g})
\qquad\text{e}\qquad
\Prob(A_{g}),
\]

la verifica che nessun regime risulti privo di osservazioni,

l'equivalenza
\[
L>\ell
\quad\Longleftrightarrow\quad
\Delta y>\frac{\ell}{D V_{0}},
\]

oppure la ricomposizione
\[
\widehat{\E}[L]
=
\sum_{g=1}^{3}
\widehat{\Prob}(A_{g})
\widehat{\E}[L\mid A_{g}].
\]

Il controllo deve essere scelto prima della scrittura del codice, perch\'{e}
esso fa parte della specifica della tappa.

\subsection{Costruire i collegamenti tra le tappe}

Una volta individuati input, operazione, output e controllo, occorre verificare
come le diverse tappe siano collegate.

Un criterio utile consiste nel chiedersi:

\begin{quote}
\small
\noindent{\ttfamily\raggedright
Quale output di questa tappa viene utilizzato successivamente?
Se questa tappa venisse eliminata, quale parte del lavoro non potrebbe essere
eseguita?
\par}
\end{quote}

Nel caso considerato, ad esempio,

\[
\text{regimi simulati}
\longrightarrow
\text{shock condizionati}
\longrightarrow
\text{perdite simulate}
\]

costituisce una dipendenza necessaria.

Le perdite simulate diventano poi input sia per l'analisi non condizionata sia
per quella condizionata:

\[
L^{(1)},\ldots,L^{(M)}
\longrightarrow
\left\{
\begin{array}{c}
\text{distribuzione complessiva},\\
\text{probabilit\`{a} di soglia},\\
\text{medie per regime}.
\end{array}
\right.
\]

La struttura del notebook assume quindi la forma di una rete ordinata di
dipendenze, anche quando viene presentata come successione lineare di celle.

\subsection{Quando due passaggi devono appartenere alla stessa tappa}

Una scomposizione eccessivamente fine produce un notebook frammentato. Una
scomposizione troppo aggregata rende invece difficile individuare errori e
controllare i risultati.

Due operazioni possono essere raggruppate nella stessa tappa quando:

\begin{itemize}
\item perseguono la stessa finalit\`{a} quantitativa;

\item utilizzano gli stessi input principali;

\item producono output strettamente collegati;

\item possono essere sottoposte allo stesso controllo logico.
\end{itemize}

Devono invece essere separate quando l'output della prima costituisce un input
necessario della seconda oppure quando richiedono controlli concettualmente
differenti.

Ad esempio, l'estrazione del regime e la generazione dello shock
condizionato possono essere considerate parti della stessa fase di simulazione,
mentre la trasformazione dello shock nella perdita appartiene a un passaggio
successivo, perch\'{e} introduce una nuova grandezza finanziaria e richiede una
verifica specifica della relazione
\[
L=D V_{0}\Delta y.
\]

\subsection{Una prima proposta di scomposizione}

Applicando i criteri precedenti al flusso teorico del caso, una prima
scomposizione pu\`{o} essere formulata come segue:

\begin{center}
\small
\begin{tabular}{p{0.9cm}p{2.4cm}p{3.0cm}p{3.0cm}p{3.0cm}}
\textbf{Tappa} &
\textbf{Input} &
\textbf{Operazione} &
\textbf{Output} &
\textbf{Controllo}
\\ \hline

1 &
Parametri finanziari e probabilistici &
Organizzare i parametri del caso &
Tabelle dei parametri &
Coerenza con la Scheda Caso
\\

2 &
Probabilit\`{a} dei regimi e parametri condizionati &
Simulare regimi e shock $\Delta y$ &
Campione dei regimi e degli shock &
Frequenze dei regimi e coerenza delle distribuzioni simulate
\\

3 &
$\Delta y$, $D$, $V_{0}$ &
Calcolare $L=D V_{0}\Delta y$ e descriverne la distribuzione &
Campione delle perdite e statistiche descrittive &
Coerenza della trasformazione e plausibilit\`{a} degli output
\\

4 &
$L$ e soglia $\ell$ &
Stimare il superamento della soglia &
$\widehat{\Prob}(L>\ell)$ &
Confronto con la soglia equivalente su $\Delta y$
\\

5 &
Perdite e regimi simulati &
Calcolare quantit\`{a} condizionate per regime &
Medie e probabilit\`{a} condizionate; stima di $\E[L\mid\G]$ &
Confronto tra gruppi e assenza di classi vuote
\\

6 &
Frequenze e medie condizionate &
Ricomporre il valore atteso complessivo &
Media ricomposta &
Verifica del valore atteso totale
\end{tabular}
\end{center}

Questa tabella costituisce una proposta iniziale, non ancora una specifica
definitiva del notebook.

\subsection{Preparare la parte variabile del Prompt 3}

La parte variabile del Prompt 3 deve mostrare la scomposizione proposta e
chiedere all'IA di verificarla senza produrre direttamente il notebook.

Una formulazione possibile \`{e}:

\begin{quote}
\small
\noindent{\ttfamily\raggedright
A partire dal flusso logico-teorico gi\`{a} validato, propongo di organizzare
il notebook in sei tappe: organizzazione dei parametri; simulazione dei regimi
e degli shock condizionati; costruzione e descrizione della perdita; analisi
del superamento della soglia; analisi condizionata per regime; verifica finale
del valore atteso totale.

Per ciascuna tappa ho individuato input, operazione, output e controllo.
Verifica se la scomposizione rispetta le dipendenze del flusso teorico, se
alcune tappe devono essere aggregate o separate e se manca qualche output o
controllo necessario. Non produrre ancora il codice Python e non modificare la
Scheda Caso.
\par}
\end{quote}

Anche in questo caso l'obiettivo non \`{e} formulare immediatamente una
struttura perfetta. L'elemento essenziale \`{e} che la prima scomposizione sia
stata proposta prima della risposta dell'IA.

\subsection{Validare la scomposizione proposta}

La risposta dell'IA pu\`{o} essere accolta soltanto dopo aver verificato che la
nuova struttura:

\begin{enumerate}
\item mantenga tutti gli elementi necessari del flusso teorico;

\item non introduca nuove ipotesi o nuove quantit\`{a} non richieste;

\item rispetti le dipendenze tra input e output;

\item associ a ogni tappa almeno un risultato osservabile;

\item preveda controlli coerenti con gli oggetti calcolati;

\item non anticipi interpretazioni che dipendono da risultati ancora non
prodotti;

\item mantenga il carattere a un periodo del modello;

\item prepari una struttura sufficientemente chiara per i successivi prompt di
tappa.
\end{enumerate}

La scomposizione validata diventa quindi la specifica operativa del notebook.

Soltanto a questo punto risulta appropriato procedere con i prompt di tappa in
Regime B e chiedere la traduzione delle singole operazioni in codice Python.

Il percorso complessivo pu\`{o} essere sintetizzato come

\[
\text{flusso teorico validato}
\longrightarrow
\text{input}
\longrightarrow
\text{operazioni}
\longrightarrow
\text{output}
\longrightarrow
\text{controlli}
\longrightarrow
\text{tappe del notebook}.
\]

Il Prompt 3, come il Prompt 2, non costituisce quindi il punto di partenza del
ragionamento. Rende esplicita e sottopone a verifica una struttura che deve
essere stata preliminarmente elaborata.

\section{Dalle tappe al notebook: simulazione, output e verifica}

Una volta validata la scomposizione in tappe input--output, la struttura
concettuale del notebook \`{e} definita. Soltanto a questo punto risulta
appropriato tradurre le singole tappe in celle Markdown, codice Python,
tabelle e grafici.

La sequenza generale \`{e}

\[
\begin{array}{c}
\text{tappa validata}
\\[3pt]
\downarrow
\\[3pt]
\text{prompt di tappa}
\\[3pt]
\downarrow
\\[3pt]
\text{risposta IA e validazione}
\\[3pt]
\downarrow
\\[3pt]
\text{cella notebook}
\\[3pt]
\downarrow
\\[3pt]
\text{output e controllo}.
\end{array}
\]

Una risposta dell'IA pu\`{o} essere sintatticamente corretta e risultare
comunque incoerente con il modello, con gli input disponibili o con gli output
richiesti. In questa fase l'IA assume quindi un ruolo operativo, mentre la
specifica teorica e la struttura delle tappe devono rimanere invariate.

\subsection{Dal prompt di tappa alla cella del notebook}

Ogni prompt di tappa parte da una struttura gi\`{a} definita. Devono essere
specificati almeno:

\begin{itemize}
\item gli input utilizzabili;

\item l'operazione da realizzare;

\item gli output richiesti;

\item il controllo da effettuare;

\item il collegamento con la tappa successiva.
\end{itemize}

Non \`{e} quindi sufficiente chiedere genericamente di ``scrivere il codice''.
Il compito dell'IA deve essere delimitato rispetto alla specifica gi\`{a}
validata.

Una formulazione tipica pu\`{o} essere:

\begin{quote}
\small
\noindent{\ttfamily\raggedright
Sviluppa la tappa indicata utilizzando esclusivamente gli input disponibili e
senza modificare la specifica teorica. Produci una breve cella Markdown, il
codice Python necessario, gli output richiesti e il controllo previsto.
Mantieni nomi delle variabili coerenti con quelli gi\`{a} utilizzati nel
notebook.
\par}
\end{quote}

La scelta di che cosa debba essere calcolato \`{e} gi\`{a} stata effettuata;
l'IA viene utilizzata principalmente per realizzare tecnicamente il calcolo.

\subsection{Il Regime B: tradurre senza modificare}

La traduzione operativa avviene prevalentemente in Regime B. Possono essere
delegate all'IA la costruzione delle strutture dati, la simulazione delle
variabili specificate, la scrittura del codice, la produzione di tabelle e
grafici e l'implementazione dei controlli gi\`{a} individuati.

La delega riguarda la realizzazione computazionale, non il significato del
problema.

Se, ad esempio,
\[
\Delta y\mid A_g
\sim
\mathcal{N}(\mu_g,\sigma_g^2),
\]
l'IA pu\`{o} proporre il codice necessario per simulare gli shock nei diversi
regimi, ma non pu\`{o} modificare distribuzione, parametri o struttura
informativa.

Analogamente, se
\[
L=D V_0\Delta y,
\]
il codice deve implementare precisamente questa trasformazione.

La regola generale \`{e}

\[
\boxed{
\text{flessibilit\`{a} tecnica}
\qquad
\text{con invarianza della specifica teorica}.
}
\]

\subsection{Simulare significa rendere osservabile il modello}

Nel caso sviluppato, la simulazione segue la struttura probabilistica del
modello. Per ogni scenario viene simulato il regime informativo
\[
Z\in\{1,2,3\},
\]
quindi lo shock
\[
\Delta y\mid A_g,
\]
e infine la perdita
\[
L=D V_0\Delta y.
\]

La dipendenza pu\`{o} essere rappresentata sinteticamente come

\[
\begin{array}{c}
Z^{(k)}
\\[3pt]
\downarrow
\\[3pt]
\Delta y^{(k)}
\\[3pt]
\downarrow
\\[3pt]
L^{(k)}.
\end{array}
\]

Il campione
\[
L^{(1)},\ldots,L^{(M)}
\]
rende osservabile la distribuzione empirica della perdita e permette di
stimare media, deviazione standard, quantili e probabilit\`{a} di superamento
di soglia.

La cella Markdown chiarisce quale oggetto viene stimato; il codice realizza il
calcolo; l'output rende osservabile il risultato.

\subsection{Un output non \`{e} ancora una verifica}

La produzione di un numero, di una tabella o di un grafico non dimostra che il
calcolo sia corretto. Ogni output deve essere confrontato con almeno una
propriet\`{a} coerente con il modello.

Per i regimi simulati deve risultare
\[
\sum_{g=1}^{3}\widehat{\Prob}(A_g)=1,
\]
con frequenze empiriche coerenti con le probabilit\`{a} teoriche.

Per la trasformazione della perdita, poich\'{e} $D>0$ e $V_0>0$,
\[
\Delta y>0
\quad\Longleftrightarrow\quad
L>0.
\]

Per una soglia $\ell>0$ deve inoltre valere
\[
L>\ell
\quad\Longleftrightarrow\quad
\Delta y>\frac{\ell}{D V_0}.
\]

Infine, nell'analisi condizionata deve essere verificata la ricomposizione
\[
\widehat{\E}[L]
=
\sum_{g=1}^{3}
\widehat{\Prob}(A_g)
\widehat{\E}[L\mid A_g].
\]

Il controllo fa quindi parte della procedura computazionale: verifica che
l'output conservi le propriet\`{a} richieste dal modello.

\subsection{Controllare anche ci\`{o} che appare plausibile}

Gli errori pi\`{u} insidiosi possono produrre risultati apparentemente
plausibili senza generare errori Python. Una frequenza simulata pu\`{o}
risultare incompatibile con la probabilit\`{a} teorica; una tabella pu\`{o}
utilizzare gruppi costruiti in modo errato; un testo pu\`{o} attribuire una
lettura dinamica a un modello a un solo periodo.

La validazione deve quindi riguardare tre livelli:

\[
\begin{array}{c}
\text{correttezza tecnica}
\\[3pt]
\downarrow
\\[3pt]
\text{coerenza quantitativa}
\\[3pt]
\downarrow
\\[3pt]
\text{coerenza con il modello finanziario}.
\end{array}
\]

La correttezza tecnica del codice non garantisce da sola la correttezza
quantitativa e interpretativa.

\subsection{Grafici: prima la domanda, poi la realizzazione}

Anche un grafico deve avere una funzione definita prima della sua costruzione.

Nel caso considerato, un istogramma della perdita pu\`{o} rendere visibili la
forma della distribuzione empirica, la dispersione e la posizione della soglia
$\ell$. Un confronto tra gruppi informativi pu\`{o} invece evidenziare
differenze tra distribuzioni o valori medi condizionati.

L'IA pu\`{o} essere delegata alla realizzazione tecnica del grafico, non alla
scelta del suo significato informativo.

La sequenza da seguire \`{e}

\[
\begin{array}{c}
\text{domanda quantitativa}
\\[3pt]
\downarrow
\\[3pt]
\text{informazione da rappresentare}
\\[3pt]
\downarrow
\\[3pt]
\text{grafico}
\\[3pt]
\downarrow
\\[3pt]
\text{codice}.
\end{array}
\]

\subsection{Quando utilizzare il Regime C}

Il Regime C viene utilizzato quando emerge una criticit\`{a} concreta. Non
\`{e} sufficiente chiedere genericamente:

\begin{quote}
\small
\noindent{\ttfamily\raggedright
Controlla se il notebook \`{e} corretto.
\par}
\end{quote}

Il prompt deve invece indicare la tappa interessata, il risultato o il codice
che genera il dubbio e il motivo per cui non appare coerente.

Ad esempio:

\begin{quote}
\small
\noindent{\ttfamily\raggedright
Nella simulazione dei regimi le frequenze empiriche risultano molto lontane
dalle probabilit\`{a} teoriche, nonostante il numero elevato di simulazioni.
Verifica il codice riportato e stabilisci se esiste una criticit\`{a}, senza
modificare i parametri teorici del caso.
\par}
\end{quote}

Un'altra criticit\`{a} pu\`{o} riguardare l'interpretazione: descrivere le
simulazioni come periodi successivi sarebbe incoerente con un modello a un
solo periodo, nel quale esse rappresentano scenari alternativi.

\subsection{Criticit\`{a} respinta e criticit\`{a} accolta}

La verifica in Regime C deve concludersi con uno dei due esiti:

\begin{center}
\begin{tabular}{p{3.5cm}p{8.0cm}}
Criticit\`{a} respinta &
il dubbio non individua un problema effettivo e il notebook non viene
modificato; \\[4pt]

Criticit\`{a} accolta &
il dubbio individua un errore o un'incoerenza e la parte interessata viene
corretta.
\end{tabular}
\end{center}

Se la criticit\`{a} viene accolta, le celle interessate devono essere
sostituite con la versione corretta. Il notebook conserva quindi il prodotto
finale validato, mentre il tracciato dell'interazione documenta il processo che
ha condotto alla correzione.

\subsection{Interpretare soltanto ci\`{o} che \`{e} stato verificato}

L'interpretazione deve seguire il controllo.

Se le medie condizionate risultano differenti nei tre regimi, \`{e} possibile
affermare che la partizione informativa distingue differenti livelli medi di
perdita nel modello simulato. Non sarebbe invece corretto attribuire tali
differenze a meccanismi economici non inclusi nella Scheda Caso oppure
interpretarle come previsioni empiriche sull'andamento futuro dei tassi.

Analogamente, una probabilit\`{a} di superamento della soglia ha significato
rispetto al modello e ai parametri assegnati, non come misura automatica del
rischio effettivo di un portafoglio reale.

Il percorso della singola tappa pu\`{o} essere sintetizzato come

\[
\begin{array}{c}
\text{specifica}
\\[3pt]
\downarrow
\\[3pt]
\text{codice}
\\[3pt]
\downarrow
\\[3pt]
\text{output}
\\[3pt]
\downarrow
\\[3pt]
\text{controllo}
\\[3pt]
\downarrow
\\[3pt]
\text{interpretazione}.
\end{array}
\]

La qualit\`{a} del notebook dipende dalla possibilit\`{a} di ricostruire
questa corrispondenza nei passaggi quantitativamente rilevanti.


\section{Dal Caso Aula al Caso TakeHome: trasferire il metodo}

Il Caso TakeHome richiede di applicare autonomamente il metodo sviluppato nelle
sezioni precedenti a un problema economico-finanziario differente.

Non si tratta di riutilizzare il notebook del Caso Aula modificando nomi e
parametri. Il nuovo caso possiede variabili, relazioni economiche, soglie e
output propri e deve quindi essere nuovamente analizzato a partire dalla
relativa Scheda Caso.

Ci\`{o} che viene trasferito \`{e} il metodo:

\[
\begin{array}{c}
\text{Scheda Caso}
\\[3pt]
\downarrow
\\[3pt]
\text{ricognizione teorica}
\\[3pt]
\downarrow
\\[3pt]
\text{flusso logico-teorico}
\\[3pt]
\downarrow
\\[3pt]
\text{tappe input--output}
\\[3pt]
\downarrow
\\[3pt]
\text{notebook}
\\[3pt]
\downarrow
\\[3pt]
\text{controlli e interpretazione}.
\end{array}
\]

L'autonomia richiesta consiste nel riconoscere questa struttura all'interno di
un problema diverso da quello gi\`{a} sviluppato.

\subsection{Un nuovo problema: lo shortfall di margine}

Il Caso TakeHome riguarda un'impresa manifatturiera energivora esposta
all'incertezza del prezzo unitario dell'energia.

La grandezza aleatoria iniziale \`{e} il prezzo $P_E$. Il costo energetico
totale e il margine operativo sono definiti rispettivamente da

\[
C_E=Q_E P_E
\]

e

\[
M=G_0-Q_E P_E.
\]

Fissato un margine minimo desiderato $m^\star$, lo shortfall di margine \`{e}

\[
S=(m^\star-M)^+.
\]

La variabile finale non \`{e} quindi una perdita lineare nello shock di tasso,
come nel Caso Aula, ma una grandezza ottenuta mediante una trasformazione che
incorpora la parte positiva.

L'informazione continua a essere rappresentata mediante una partizione

\[
\Omega=A_1\cup A_2\cup A_3,
\qquad
A_i\cap A_j=\emptyset
\quad\text{per }i\neq j,
\]

con

\[
P_E\mid A_g
\sim
\mathcal{N}(\mu_g,\sigma_g^2).
\]

I tre eventi rappresentano normalizzazione energetica, tensione persistente e
shock energetico severo. Rimane quindi centrale il collegamento tra
informazione, distribuzioni condizionate, probabilit\`{a} di soglia e valori
attesi condizionati, mentre cambia la struttura economica che conduce alla
variabile finale.

\subsection{Che cosa rimane invariato}

Tra Caso Aula e Caso TakeHome rimangono invariati alcuni principi:

\begin{itemize}
\item partire dalla Scheda Caso senza modificarne la specifica;

\item riconoscere la variabile finale e le grandezze da cui dipende;

\item individuare il ruolo della partizione informativa;

\item distinguere quantit\`{a} teoriche e stime Monte Carlo;

\item costruire il flusso logico-teorico prima del notebook;

\item trasformarlo in tappe input--output con controlli espliciti;

\item utilizzare l'IA nelle fasi di ricognizione, traduzione operativa e
verifica;

\item interpretare i risultati entro i limiti del modello.
\end{itemize}

Rimangono inoltre centrali

\[
\E[S\mid A_g],
\]

\[
\E[S\mid\G]
=
\sum_{g=1}^{3}
\E[S\mid A_g]\1_{A_g},
\]

e

\[
\E[S]
=
\sum_{g=1}^{3}
\Prob(A_g)\E[S\mid A_g].
\]

La struttura probabilistica del condizionamento rimane quindi invariata,
mentre cambia la grandezza economica alla quale viene applicata.

\subsection{Che cosa deve essere ricostruito}

Il trasferimento del metodo non deve diventare una copia meccanica.

Deve innanzitutto essere ricostruita la catena economica

\[
P_E
\longrightarrow
C_E
\longrightarrow
M
\longrightarrow
S.
\]

Devono inoltre essere riconosciute le soglie implicite. Dalla definizione dello
shortfall segue

\[
S>0
\quad\Longleftrightarrow\quad
P_E>
\frac{G_0-m^\star}{Q_E},
\]

mentre, per una soglia severa $s^\star$,

\[
S>s^\star
\quad\Longleftrightarrow\quad
P_E>
\frac{G_0-m^\star+s^\star}{Q_E}.
\]

Queste equivalenze collegano il fattore di rischio alle condizioni economiche
di interesse e costituiscono anche controlli del notebook.

Deve inoltre essere riconosciuta la particolare forma di

\[
S=(m^\star-M)^+.
\]

Per costruzione, lo shortfall non pu\`{o} essere negativo e pu\`{o} presentare
una massa di probabilit\`{a} in zero.

Infine, la distribuzione normale assegnata a $P_E$ pu\`{o} produrre prezzi
negativi. Essi non devono essere eliminati automaticamente, ma rilevati e
discussi come limite della specificazione prevista dalla Scheda Caso.

\subsection{Ricostruire autonomamente il Prompt 2}

Nel Caso TakeHome la parte variabile del Prompt 2 deve essere nuovamente
preparata prima dell'interazione con l'IA.

Le domande da utilizzare rimangono quelle introdotte nel Caso Aula: quale sia
la variabile finale, da quali grandezze dipenda, quale sia la fonte
dell'incertezza, quale ruolo svolga l'informazione, quali quantit\`{a} debbano
essere calcolate e quali propriet\`{a} possano diventare controlli.

Le risposte sono per\`{o} specifiche del nuovo problema. La proposta iniziale
dovr\`{a} quindi riconoscere almeno la partizione informativa, la distribuzione
condizionata di $P_E$, la catena che conduce a $S$, le probabilit\`{a} di
soglia, i valori attesi condizionati e la verifica del valore atteso totale.

Non \`{e} necessario che l'ordine sia gi\`{a} definitivo, ma deve essere
visibile una ricognizione teorica autonoma.

\subsection{Ricostruire autonomamente il Prompt 3}

Dopo la validazione del nuovo flusso teorico deve essere costruita una nuova
scomposizione operativa. Non \`{e} necessario che numero e contenuto delle
tappe coincidano con quelli del Caso Aula.

Per ogni tappa devono essere nuovamente definiti input, operazione, output e
controllo.

Ad esempio, la costruzione dello shortfall richiede come input il margine
simulato e $m^\star$, applica

\[
S=(m^\star-M)^+,
\]

produce il vettore degli shortfall e deve verificare almeno

\[
S\geq0.
\]

Analogamente, l'analisi delle soglie deve confrontare gli eventi definiti su
$S$ con le corrispondenti soglie implicite su $P_E$.

Il Prompt 3 deve quindi utilizzare la stessa logica di decomposizione appresa
nel Caso Aula, ma applicata alla struttura specifica del nuovo problema.

\subsection{Trasferire il metodo, non la soluzione}

Il passaggio dal Caso Aula al Caso TakeHome pu\`{o} essere sintetizzato nella
distinzione seguente:

\begin{center}
\small
\begin{tabular}{p{5.2cm}p{6.6cm}}
\textbf{Si trasferisce} & \textbf{Deve essere ricostruito} \\
\hline
metodo di lettura della Scheda Caso &
catena economica delle variabili \\

procedura per costruire il Prompt 2 &
flusso logico-teorico specifico \\

schema input--operazione--output--controllo &
tappe effettive del notebook \\

distinzione tra teoria e stima empirica &
quantit\`{a} specifiche da stimare \\

uso dei Regimi A, B e C &
prompt concretamente necessari \\

principio di validazione &
controlli adatti al nuovo modello
\end{tabular}
\end{center}

Affrontare autonomamente un nuovo caso non significa quindi riprodurre una
soluzione gi\`{a} osservata. Significa distinguere gli elementi generali del
metodo da quelli che devono essere ricostruiti a partire dalla nuova specifica.
Il Caso TakeHome verifica precisamente questa capacit\`{a} di trasferimento.

\section{Sintesi finale}

Questo capitolo ha mostrato come gli strumenti teorici introdotti nei Capitoli
1--3 possano essere utilizzati nella soluzione di casi quantitativi mediante
Python e intelligenza artificiale.

Eventi, partizioni informative, variabili casuali, distribuzioni,
probabilit\`{a} di soglia, valori attesi e valori attesi condizionati non
vengono sostituiti dalla simulazione. Vengono tradotti in quantit\`{a}
empiriche che possono essere calcolate, rappresentate e controllate nel
notebook.

La sequenza generale pu\`{o} essere sintetizzata come

\[
\begin{array}{c}
\text{Scheda Caso}
\\[3pt]
\downarrow
\\[3pt]
\text{flusso logico-teorico}
\\[3pt]
\downarrow
\\[3pt]
\text{tappe input--output}
\\[3pt]
\downarrow
\\[3pt]
\text{notebook}
\\[3pt]
\downarrow
\\[3pt]
\text{controlli e interpretazione}.
\end{array}
\]

I passaggi pi\`{u} delicati precedono la scrittura del codice.

Il \textbf{Prompt 2} richiede di ricostruire il flusso logico-teorico
risolutivo. La sua parte variabile deve rendere riconoscibili almeno la
variabile finale di interesse, le grandezze da cui essa dipende, la fonte
dell'incertezza, il ruolo dell'informazione, le quantit\`{a} da calcolare e le
propriet\`{a} teoriche utilizzabili come collegamenti o controlli.

Il Prompt 2 non serve quindi a chiedere all'IA di individuare autonomamente la
soluzione, ma a sottoporre a verifica una prima struttura teorica gi\`{a}
riconosciuta.

Il \textbf{Prompt 3} interviene dopo la validazione del flusso e trasforma le
dipendenze teoriche in una struttura computazionale. Ogni tappa deve essere
riconducibile allo schema

\[
\text{input}
\longrightarrow
\text{operazione}
\longrightarrow
\text{output}
\longrightarrow
\text{controllo}.
\]

Anche in questo caso deve esistere una proposta iniziale. L'IA pu\`{o} aiutare
a verificare l'ordine delle tappe, i collegamenti tra input e output e
l'adeguatezza dei controlli.

Soltanto dopo questi passaggi risulta appropriato utilizzare l'IA per la
traduzione operativa in codice. Python rende computabile una specifica teorica
gi\`{a} definita, mentre l'IA pu\`{o} contribuire alla realizzazione tecnica
delle celle, delle tabelle, dei grafici e dei controlli.

La risposta dell'IA non deve essere trasferita automaticamente nel notebook.
Il processo rimane

\[
\begin{array}{c}
\text{prompt}
\\[3pt]
\downarrow
\\[3pt]
\text{risposta IA}
\\[3pt]
\downarrow
\\[3pt]
\text{validazione}
\\[3pt]
\downarrow
\\[3pt]
\text{notebook}.
\end{array}
\]

La validazione riguarda non soltanto l'esecuzione tecnica del codice, ma anche
la coerenza quantitativa degli output e la compatibilit\`{a} con il modello
finanziario assegnato.

Quando emerge un dubbio specifico, il Regime C consente di utilizzare l'IA
come strumento di verifica critica. La criticit\`{a} deve essere esplicitamente
respinta oppure accolta; nel secondo caso viene corretto il notebook, mentre il
processo di verifica rimane documentato nel tracciato dell'interazione.

Il confronto tra il caso sviluppato e il Caso TakeHome mostra infine che deve
essere trasferito il metodo, non una soluzione specifica. Cambiano variabili,
relazioni economiche, soglie e quantit\`{a} di interesse; rimangono invariati i
principi con cui il problema viene analizzato, scomposto, reso computabile e
verificato.

L'obiettivo finale consiste quindi nel rendere progressivamente autonomo il
passaggio

\[
\begin{array}{c}
\text{teoria}
\\[3pt]
\downarrow
\\[3pt]
\text{struttura della soluzione}
\\[3pt]
\downarrow
\\[3pt]
\text{procedura computazionale}
\\[3pt]
\downarrow
\\[3pt]
\text{risultati controllati}.
\end{array}
\]

Python e l'intelligenza artificiale svolgono funzioni importanti all'interno
di questo percorso, ma il significato del problema, la scelta degli oggetti
teorici rilevanti e la valutazione della coerenza dei risultati rimangono
elementi essenziali della soluzione quantitativa.